\uuid{12z4}
\exo7id{1189}
\titre{exo7 1189}
\auteur{barraud}
\organisation{exo7}
\datecreate{2003-09-01}
\isIndication{false}
\isCorrection{false}
\chapitre{Déterminant, système linéaire}
\sousChapitre{Système linéaire, rang}
\cpge{mpsi-mp2i.espaces-vectoriels.b, pcsi.espaces-vectoriels.b, ptsi.espaces-vectoriels.a, mpsi-mp2i.matrices-applications-lineaires.c, pcsi.matrices-applications-lineaires.d, ptsi.matrices-applications-lineaires.d, bcpst1.espace-vectoriel-kn.a, bcpst1.systemes-lineaires.a, bcpst2.espaces-vectoriels.a}
\module{Algèbre}
\niveau{L2}
\difficulte{}

\contenu{
\question{
Soient $v_1=\begin{pmatrix} a\\b\\c \end{pmatrix}$ et $v_2=\begin{pmatrix} a'\\b'\\c' \end{pmatrix}$ deux
  vecteurs indépendants de $\R^3$. Donner, sous forme d'équation, une
  condition nécessaire et suffisante pour qu'un vecteur
  $w=\begin{pmatrix} x\\y\\z \end{pmatrix}$ appartienne à l'espace vectoriel engendré par
  $v_1$ et $v_2$.
  
  Même question pour un plan engendré par deux vecteurs de $\R^4$.
}
}
